%% kapa: exam=1 round=33 year=2022 period=1 form=A subject=economics subject_ko=경제학원론 seq=25 number=65
%% source: Q-Net 감정평가사 기출문제 (cst00309) artlSeq=5209997 file=raw/1008/5209997/2022년도 제33회 감정평가사 1차시험 1교시(민법,경제학,부동산학).pdf
%% conversion: kordoc 4.24.51 -> tools/gen1.py (auto); review=proofread-sonnet; flags=none
\begin{question}{65}
A국 중앙은행은 아래의 테일러 규칙(Taylor rule)에 따라 명목정책금리를 조정한다. 이에 관한 설명으로 옳지 않은 것은? \{단, 총생산 갭 = (실질GDP $-$ 완전고용 실질GDP) / 완전고용 실질GDP이다.\}
\begin{framedtext}
명목정책금리 = 인플레이션율 + 0.02 + 0.5 $\times$ (인플레이션율 $-$ 0.03) + 0.5 $\times$ (총생산 갭)
\end{framedtext}
\begin{choices}
  \item A국 중앙은행의 인플레이션율 목표치는 3\%이다.
  \item 인플레이션율 목표치를 2\%로 낮추려면 명목정책금리를 0.5\%p 인하해야 한다.
  \item 인플레이션율이 목표치와 동일하고 총생산 갭이 1\%인 경우 실질 이자율은 2.5\%이다.
  \item 완전고용 상태에서 인플레이션율이 2\%인 경우에 명목정책금리는 3.5\%로 설정해야 한다.
  \item 인플레이션율이 목표치보다 1\%p 더 높은 경우에 명목정책금리를 0.5\%p 인상한다.
\end{choices}
\end{question}
\answer{2}
